\begin{answers}{ch:capstone}{総仕上げ：自作 LLM を最初から最後まで}
\item GPU 時間は約 7.6、日数は約 0.08 日（約 1.9 時間）、費用は約 30 ドルです。
計算量は $C=6ND=6\times(3\times10^{8})\times(6\times10^{9})=1.08\times10^{19}$ FLOPs です。
H100 1 枚が 1 時間にこなす計算は、$989\times10^{12}\times0.4\times3600\approx1.42\times10^{18}$ FLOPs です。
割ると $1.08\times10^{19}/(1.42\times10^{18})\approx7.6$ GPU 時間になります。
4 枚で分けるので $7.6/4\approx1.9$ 時間、24 で割って約 0.08 日です。
費用は $7.6\times3.99\approx30$ ドルです。
MFU が 25\% なら、GPU 時間は $0.4/0.25=1.6$ 倍の約 12.1 に増えます。
日数は約 0.13 日（約 3.0 時間）、費用は約 48 ドルです。
MFU は式の分母にあるので、日数も費用も MFU に反比例します。
\item 結果は用意したテキストで変わるので、見通しを書きます（実測ではありません）。
字/トークンは、まず文字の種類で決まります。
BPE はバイトから出発し、半角の英数字は 1 字 1 バイト、かなや漢字は 1 字 3 バイトだからです。
本文の縮小版でも、計算ドリルの行は約 1.6 字/トークン、漱石の行は約 0.8 字/トークンでした。
英語の文章や、同じ言い回しをくり返す文章（日誌など）なら、漱石より大きくなるはずです。
検証損失は 1 トークンあたりの値なので、字/トークンで割って 1 字あたりに直してから比べます。
もう一つ、\code{mini} は渡したテキストに計算ドリルの行を混ぜるので、検証損失はテキストの部分とドリルの部分の平均です。
漱石のときの 3.405 は、小説の部分 3.62 とドリルの部分 2.77 を、トークン数で 72 対 28 に混ぜた値でした（\secref{sec:cap-mini}）。
テキストの長さが違えば混ざり方も変わるので、比べるのは、テキストの部分だけの損失です（\file{ch25/sweep_scale.py} と同じく、検証データを分けて測ります）。
言い回しが限られた文章ほど次のトークンを当てやすく、損失は小さくなると予想できます。
\item 所要時間は約 2.7 倍、計算量 $6ND$ は約 3.1 倍になりました（実測。乱数の種は 1 つ）。
パラメータ数を\eqref{eq:tf-params}で数えると、$512\times128+64\times128+4\times(12\times128^{2}+4\times128)+2\times128=862{,}464$ です。
元の 277,440 の約 3.1 倍で、$D$ は変わらないので、$6ND$ も約 3.1 倍です。
同じ計算機で続けて測ると、元の設定は 84 秒、新しい設定は 225 秒でした（本文の 110 秒とは混み具合が違います）。
検証損失は 3.405 から 3.151 に下がりました。
しかし対話の正答率は、事前学習後 0.0\%、SFT 後 10.0\%、強化学習後 10.0\%、int8 で 10.0\% でした。
元の設定（0.0\%、11.0\%、10.0\%、10.0\%）との差は、200 問の標準誤差（約 2.1 ポイント）より小さく、区別できません。
\item 掛け算では 100 問すべてで電卓を呼び、加減算の 100 問では 1 度も呼びませんでした。
しかし、掛け算の正答率は 23\% でした（実測。乱数の種は 1 つ）。
SFT の段の手本を、加減算 2,000 件と、\code{make_tool_example} の \code{mul} で作った掛け算 2,000 件に替えました。
どちらも \code{encode_chat_example} に \code{learn_role=True} を渡します。
役割トークンも正解に含めないと、応答を \code{<|tool|>} と \code{<|assistant|>} のどちらで始めるかを学べないからです。
評価では、第 23 章の \code{run_agent} に縮小版のモデルとトークナイザをつなぎました。
電卓に渡した式が正しかったのは 58\% です（「11と57を掛けると」に \code{11*55}）。
結果を応答へ写すところでも間違えます（結果 \code{605} に「630です。」）。
2 層 96 次元のモデルは、「いつ呼ぶか」は覚えても、数字を正しく写すことがまだできません。
\item 例として案 (3)、電卓を持つエージェントを選びます。
ふつうの問題、境目の問題、答えられない問題を混ぜます。
電卓なしで答える問題は「12+35は？」「999+1は？」「1000-1は？」「50から3を引くといくつですか？」です。
電卓を呼ぶ問題は「37×48は？」「0×512は？」「7と13を掛けるといくつですか？」「123×456+789は？」「999×999は？」です。
最後に「今日の天気は？」（できないと答える）を入れます。
正解は 47、1000、999、47、1776、0、91、56877、998001 と、プログラムで照合できる形に決めます。
成功の基準は測る前に決めます。
たとえば「公開モデルをそのまま使ったときとの差が、差の標準誤差の 2 倍を超える」です。
二つを別々の問題で測るなら、差の標準誤差は、1 つの標準誤差の $\sqrt2$ 倍です（同じ問題で測れば、対応ありの式\eqref{eq:es-paired}で、もう少し小さくなりえます）。
1 つの標準誤差が 3 ポイントなら、差の標準誤差は約 4.2 ポイントで、成功と言えるのは約 8.5 ポイントより大きい差です。
標準誤差は $p(1-p)/n$ の平方根なので、0.03 以下にするには $n$ が $p(1-p)/0.03^{2}$ 以上でなければなりません。
いちばん厳しい $p=0.5$ で $0.25/0.0009\approx278$ 問、$p=0.9$ なら $0.09/0.0009=100$ 問です。
10 問では標準誤差が約 16 ポイントあるので、10 問は型を決める見本とし、プログラムで数百問に増やします。
\item 例として、第 22 章の「能力を伸ばす強化学習は伸びが小さい」を選びます。
260 万パラメータのモデルでは、伸びは 2.1 ポイント（82.6\% → 84.7\%）でした。
1 つのモデルの標準誤差は約 1.2 ポイントなので、二つのモデルの差の標準誤差は、別々の問題で測ったとすれば $\sqrt2\times1.2\approx1.7$ ポイントです（\chapref{ch:rl-llm}、\secref{sec:cap-map}）。
2.1 ポイントはその約 1.2 倍で、2 倍の約 3.4 ポイントに届かず、偶然の差と区別しきれません。
同じ 1,000 問で比べているので、対応ありの式\eqref{eq:es-paired}を使えば差の標準誤差はもう少し小さくなりえますが、本書ではそれで判定していません。
変わりうる理由の一つは、強化学習が強められるのは、サンプリングでたまたま出た正解だけだということです。
正答率が 1 割ほどの縮小版では、8 本とも不正解のグループが 4 割ほどあり、そのグループは重みを動かしませんでした（縮小版の強化学習が伸びなかったのは、学習の量が少なく、200 問の評価が粗いことも理由です）。
大きなモデルが難しい問題でも何本かに 1 本は正解を出せるなら、学ぶ手がかりが増えます。
第 24 章で強化学習が大きく効いたのも、正解に届く行動（電卓の正しい呼び方）を、学習の前から身に付けていたからでした（\secref{sec:cap-map}）。
確かめるには、層と次元を変えたモデルを 3〜4 段用意し、同じデータで事前学習と SFT をします（\figref{fig:cap-sweep}のスイープを、強化学習の段まで延ばす形です）。
そこへ同じ設定の GRPO をかけ、固定の 1,000 問で前後の正答率を測ります（乱数の種は 3 つ以上）。
各段で、伸びが差の標準誤差（各段の正答率から求めた値。82〜85\% 付近なら約 1.7 ポイント）の 2 倍を超えるか、規模とともに伸びが大きくなるかを見ます。
比べる相手として、正解のお手本で SFT を続けたモデルも測ります（第 22 章では、こちらが 88.6\% で強化学習より伸びました）。
これは予想で、本書では確かめていません。
\end{answers}
